\documentclass[11pt]{article}
\usepackage{amsmath, amssymb, amsthm}
\usepackage[margin=1in]{geometry}
\usepackage{enumitem}
\usepackage{booktabs}
\usepackage{algorithm}
\usepackage{algpseudocode}

\title{Approach D: Guided Reverse Diffusion for\\Utility-Constrained Image Generation}
\author{Discussion Notes}
\date{March 2026}

\newcommand{\xb}{\mathbf{x}}
\newcommand{\xc}{\xb_{\mathrm{c}}}
\newcommand{\zb}{\mathbf{z}}
\newcommand{\eb}{\boldsymbol{\epsilon}}
\newcommand{\gb}{\mathbf{g}}
\newcommand{\Ib}{\mathbf{I}}
\newcommand{\R}{\mathbb{R}}
\newcommand{\E}{\mathbb{E}}
\newcommand{\alphabar}{\bar{\alpha}}

\begin{document}
\maketitle

\tableofcontents
\vspace{1em}

% ============================================================================
\section{Goal}
% ============================================================================

We want to generate an image $\xb_0 \in \R^d$ (where $d = 64 \times 64 = 4096$)
that simultaneously:
\begin{enumerate}[nosep]
    \item Looks like a natural image (lies on or near the natural image manifold $\mathcal{M}$).
    \item Has high GP-based distribution-aware utility $U(\xb_0)$ for a target neuron.
\end{enumerate}

Approaches A--C (from \texttt{GUIDANCE\_POSSIBILITIES.tex}) all start from an existing image
and optimize it, using the diffusion model as a regularizer \emph{after the fact}.
Their common weakness: the diffusion model's score is evaluated in situations it wasn't
trained for (clean images, amplified images, etc.).

Approach D flips the paradigm: we \textbf{generate} the image through the diffusion model's
reverse process, nudging each denoising step toward high utility. The diffusion model
is always in-distribution (it processes images at exactly the noise level it expects),
and the utility acts as a guidance signal.

% ============================================================================
\section{Prerequisites and Notation}
% ============================================================================

\subsection{The trained DDPM}

We have a trained noise predictor $\eb_\theta(\xb_t, t)$ (our UNet, 2.16M parameters)
and a cosine noise schedule with:
\begin{itemize}[nosep]
    \item $T = 1000$ timesteps
    \item $\alphabar_t = \prod_{s=1}^t (1 - \beta_s)$ --- cumulative signal retention
    \item $\beta_t = 1 - \alphabar_t / \alphabar_{t-1}$ --- noise added at step $t$
    \item $\alpha_t = 1 - \beta_t$ --- signal retained at step $t$
\end{itemize}

Key boundary values from our cosine schedule:
\begin{center}
\begin{tabular}{lcccc}
\toprule
$t$ & $\alphabar_t$ & $\sqrt{\alphabar_t}$ & $\sqrt{1-\alphabar_t}$ & $\beta_t$ \\
\midrule
1   & 0.9999 & 1.000  & 0.010  & 0.0000 \\
50  & 0.9939 & 0.997  & 0.078  & 0.0001 \\
250 & 0.8445 & 0.919  & 0.394  & 0.0011 \\
500 & 0.4938 & 0.703  & 0.711  & 0.0031 \\
750 & 0.1310 & 0.362  & 0.930  & 0.0070 \\
999 & 0.000002 & 0.001 & 1.000 & 0.7500 \\
\bottomrule
\end{tabular}
\end{center}

\subsection{Key identities}

\paragraph{Forward process.}
Given a clean image $\xb_0$, the noisy version at time $t$ is:
\begin{equation}
    \xb_t = \sqrt{\alphabar_t}\, \xb_0 + \sqrt{1 - \alphabar_t}\, \eb,
    \qquad \eb \sim \mathcal{N}(\mathbf{0}, \Ib)
    \label{eq:forward}
\end{equation}

\paragraph{Tweedie's formula.}
Given $\xb_t$ at noise level $t$, the posterior mean of $\xb_0$ is:
\begin{equation}
    \hat{\xb}_0(\xb_t, t)
    = \frac{\xb_t - \sqrt{1 - \alphabar_t}\, \eb_\theta(\xb_t, t)}{\sqrt{\alphabar_t}}
    \label{eq:tweedie}
\end{equation}

\paragraph{Score function.}
The score of the noisy distribution $p_t(\xb_t)$ is:
\begin{equation}
    \nabla_{\xb_t} \log p_t(\xb_t)
    \approx s_\theta(\xb_t, t)
    = -\frac{\eb_\theta(\xb_t, t)}{\sqrt{1 - \alphabar_t}}
    \label{eq:score}
\end{equation}

\subsection{The GP utility}

$U(\xb) = H_{\text{marg}}(\xb) - H_{\text{cond}}(\xb)$ is the distribution-aware
utility, computed via \texttt{distribution\_aware\_utility()} from \texttt{acquisition.py}.
It measures information gain about the target neuron's response.

$U$ is differentiable with respect to $\xb$ through the GP kernel (arc-cosine kernel
with RF structure). The gradient $\nabla_\xb U(\xb)$ is well-defined for
images in the valid pixel range.

We also have a firing rate guard: if the log-firing rate $\mu_g(\xb) = A \cdot \lambda(\xb) + \lambda_0$
exceeds $\log(f_{\max})$, the utility is considered invalid (physiologically unrealistic firing rate).

% ============================================================================
\section{Standard DDPM Reverse Process (Unconditional)}
% ============================================================================

The standard (unguided) reverse process generates images by iterating:
\begin{equation}
    \xb_{t-1} = \frac{1}{\sqrt{\alpha_t}}
    \left(\xb_t - \frac{\beta_t}{\sqrt{1 - \alphabar_t}}\, \eb_\theta(\xb_t, t)\right)
    + \sigma_t\, \zb
    \label{eq:ddpm_reverse}
\end{equation}
where $\zb \sim \mathcal{N}(\mathbf{0}, \Ib)$ for $t > 1$ and $\zb = \mathbf{0}$ for $t = 1$,
and $\sigma_t = \sqrt{\beta_t}$.

Starting from $\xb_T \sim \mathcal{N}(\mathbf{0}, \Ib)$, after $T$ steps we obtain
$\xb_0 \sim p_\theta(\xb_0) \approx p_{\text{data}}(\xb_0)$.

Equivalently, using the Tweedie estimate $\hat{\xb}_0$ from Eq.~\eqref{eq:tweedie}:
\begin{equation}
    \xb_{t-1} = \sqrt{\alphabar_{t-1}}\, \hat{\xb}_0
    + \sqrt{1 - \alphabar_{t-1} - \sigma_t^2}\,
    \frac{\xb_t - \sqrt{\alphabar_t}\, \hat{\xb}_0}{\sqrt{1 - \alphabar_t}}
    + \sigma_t\, \zb
    \label{eq:ddpm_reverse_x0}
\end{equation}

This ``predicted $\xb_0$'' form will be more natural for adding guidance.

% ============================================================================
\section{Classifier Guidance: Theory}
\label{sec:classifier_guidance}
% ============================================================================

\subsection{The conditional score}

We want to sample from $p(\xb_0 \mid y)$ where $y$ represents ``high utility.''
Modeling this as $p(y \mid \xb_0) \propto \exp(w \cdot U(\xb_0))$ for guidance scale $w > 0$,
Bayes' rule gives:
\begin{equation}
    \nabla_{\xb_t} \log p(\xb_t \mid y)
    = \underbrace{\nabla_{\xb_t} \log p_t(\xb_t)}_{\text{unconditional score}}
    + \underbrace{\nabla_{\xb_t} \log p(y \mid \xb_t)}_{\text{classifier gradient}}
    \label{eq:conditional_score}
\end{equation}

The unconditional score is provided by the trained UNet via Eq.~\eqref{eq:score}.
The classifier gradient requires evaluating $p(y \mid \xb_t)$ --- but $\xb_t$ is noisy,
and our utility $U$ is defined on clean images.

\subsection{Approximation: utility of the Tweedie estimate}

The standard approach (Dhariwal \& Nichol, 2021) evaluates the classifier on
$\hat{\xb}_0(\xb_t, t)$ instead of the noisy $\xb_t$:
\begin{equation}
    \nabla_{\xb_t} \log p(y \mid \xb_t)
    \approx w \cdot \nabla_{\xb_t} U\!\left(\hat{\xb}_0(\xb_t, t)\right)
    \label{eq:guidance_gradient}
\end{equation}

This gradient flows through two stages:
\begin{enumerate}[nosep]
    \item $\nabla_{\hat{\xb}_0} U(\hat{\xb}_0)$ --- gradient of utility w.r.t.\ the clean image estimate
    \item $\frac{\partial \hat{\xb}_0}{\partial \xb_t}$ --- Jacobian of the Tweedie formula w.r.t.\ the noisy image
\end{enumerate}

We define the \textbf{guidance gradient}:
\begin{equation}
    \gb_t \;=\; \nabla_{\xb_t} U\!\left(\hat{\xb}_0(\xb_t, t)\right)
    \label{eq:g_def}
\end{equation}

\subsection{Chain rule through Tweedie}

From Eq.~\eqref{eq:tweedie}:
\begin{equation}
    \hat{\xb}_0 = \frac{1}{\sqrt{\alphabar_t}}
    \left(\xb_t - \sqrt{1 - \alphabar_t}\, \eb_\theta(\xb_t, t)\right)
\end{equation}

The Jacobian $\frac{\partial \hat{\xb}_0}{\partial \xb_t}$ has two contributions:
\begin{equation}
    \frac{\partial \hat{\xb}_0}{\partial \xb_t}
    = \frac{1}{\sqrt{\alphabar_t}}
    \left(\Ib - \sqrt{1 - \alphabar_t}\, \frac{\partial \eb_\theta}{\partial \xb_t}\right)
    \label{eq:tweedie_jacobian}
\end{equation}

where $\frac{\partial \eb_\theta}{\partial \xb_t} \in \R^{d \times d}$ is the Jacobian
of the UNet. This is a $d \times d$ matrix ($d = 4096$), which we never form explicitly.
Instead, PyTorch's autograd computes the vector-Jacobian product (VJP) directly:
\begin{equation}
    \gb_t = \left(\nabla_{\hat{\xb}_0} U\right)^T
    \frac{\partial \hat{\xb}_0}{\partial \xb_t}
\end{equation}

This is computed automatically by calling \texttt{loss.backward()} where
$\texttt{loss} = -U(\hat{\xb}_0)$ and $\xb_t$ is a leaf tensor with \texttt{requires\_grad=True}.

\subsection{Modified noise prediction}

The guided reverse step modifies the noise prediction to incorporate the classifier:
\begin{equation}
    \eb_{\text{guided}} = \eb_\theta(\xb_t, t)
    - w\,\sqrt{1 - \alphabar_t}\, \gb_t
    \label{eq:eps_guided}
\end{equation}

\textbf{Derivation.} The unconditional score is $s = -\eb_\theta / \sqrt{1-\alphabar_t}$.
The conditional score is $s + w \cdot \gb_t$. Converting back to noise prediction:
\begin{align}
    \eb_{\text{guided}}
    &= -\sqrt{1-\alphabar_t}\,(s + w \cdot \gb_t) \\
    &= \eb_\theta - w\,\sqrt{1-\alphabar_t}\, \gb_t
\end{align}

The guided Tweedie estimate is:
\begin{equation}
    \hat{\xb}_0^{\text{guided}}
    = \frac{\xb_t - \sqrt{1-\alphabar_t}\, \eb_{\text{guided}}}{\sqrt{\alphabar_t}}
    = \hat{\xb}_0 + \frac{w\,(1-\alphabar_t)}{\sqrt{\alphabar_t}}\, \gb_t
    \label{eq:x0_guided}
\end{equation}

The second term shows that guidance shifts the clean image estimate in the direction
of $\gb_t$ (increasing utility), with a weight that grows with the noise level
$(1 - \alphabar_t)$. At $t=1$, the shift is negligible ($\approx 10^{-4}$).
At $t=500$, it is $\approx 0.7 w$.

% ============================================================================
\section{Full Algorithm: DDPM with Classifier Guidance}
\label{sec:algorithm_ddpm}
% ============================================================================

\begin{algorithm}[H]
\caption{Guided DDPM Reverse Process}
\label{alg:guided_ddpm}
\begin{algorithmic}[1]
\Require Trained UNet $\eb_\theta$, schedule $\{\alphabar_t, \beta_t, \alpha_t\}$,
         GP model + likelihood, guidance scale $w$, firing rate cap $f_{\max}$
\Ensure Generated image $\xb_0$ with high utility
\State $\xb_{T-1} \sim \mathcal{N}(\mathbf{0}, \Ib)$
    \Comment{Start from noise (skip $t=T$ per cosine schedule fix)}
\For{$t = T{-}1, T{-}2, \ldots, 1$}
    \State \textbf{--- UNet forward (with gradient tracking) ---}
    \State $\xb_t.\text{requires\_grad} \gets \text{True}$
    \State $\hat{\eb} \gets \eb_\theta(\xb_t, t)$
        \Comment{UNet forward pass}
    \State $\hat{\xb}_0 \gets (\xb_t - \sqrt{1-\alphabar_t}\,\hat{\eb}) / \sqrt{\alphabar_t}$
        \Comment{Tweedie estimate}
    \Statex
    \State \textbf{--- Utility evaluation on Tweedie estimate ---}
    \State Convert $\hat{\xb}_0$ from diffusion space to raw pixel space
    \State $U, \mu_g \gets \texttt{distribution\_aware\_utility}(\hat{\xb}_0)$
    \Statex
    \State \textbf{--- Compute guidance gradient ---}
    \If{$\exp(\mu_g) < f_{\max}$}
        \Comment{Only guide if firing rate is physical}
        \State $(-U).\text{backward}()$
            \Comment{Backprop through Tweedie + UNet}
        \State $\gb_t \gets -\xb_t.\text{grad}$
            \Comment{$\gb_t = \nabla_{\xb_t} U(\hat{\xb}_0)$}
    \Else
        \State $\gb_t \gets \mathbf{0}$
    \EndIf
    \Statex
    \State \textbf{--- Guided reverse step ---}
    \State $\hat{\eb}_{\text{guided}} \gets \hat{\eb} - w\sqrt{1-\alphabar_t}\,\gb_t$
    \State $\mu_t \gets \frac{1}{\sqrt{\alpha_t}}(\xb_t - \frac{\beta_t}{\sqrt{1-\alphabar_t}}\hat{\eb}_{\text{guided}})$
    \If{$t > 1$}
        \State $\xb_{t-1} \gets \mu_t + \sqrt{\beta_t}\,\zb$,
            \quad $\zb \sim \mathcal{N}(\mathbf{0}, \Ib)$
    \Else
        \State $\xb_0 \gets \mu_t$
    \EndIf
    \State Detach $\xb_{t-1}$ from computation graph
\EndFor
\State \Return $\xb_0$
\end{algorithmic}
\end{algorithm}

\subsection{Important details}

\paragraph{Detach after each step.}
After computing $\xb_{t-1}$, we detach it from the computation graph.
Each step is an independent gradient computation --- we do \emph{not} backpropagate
through the entire chain of $T$ steps (that would require $O(T)$ memory for all
intermediate UNet activations).

\paragraph{Skip $t = T$.}
Our cosine schedule has $\alpha_{1000} \approx 0.001$, giving $1/\sqrt{\alpha_{1000}} \approx 31.6$.
Starting the loop at $t = T{-}1 = 999$ avoids this amplification (same fix as unconditional sampling).

\paragraph{Guidance gradient computation.}
At each step, we compute $\gb_t$ by:
\begin{enumerate}[nosep]
    \item Setting $\xb_t$ to require gradients
    \item Running the UNet forward (gradients enabled, but UNet weights frozen)
    \item Computing $\hat{\xb}_0$ via Tweedie
    \item Evaluating $U(\hat{\xb}_0)$ through the GP
    \item Calling \texttt{backward()} to get $\xb_t.\text{grad}$
\end{enumerate}

The UNet weights are frozen (\texttt{unet.requires\_grad\_(False)}), so no optimizer
updates to $\theta$ occur. The only purpose of the backward pass is to compute
$\partial U / \partial \xb_t$.

\paragraph{Space conversions.}
The UNet operates in diffusion space (approximately $[-1, 1]$, normalized by \texttt{scale\_factor} $= 2.478$).
The GP operates in raw pixel space. The conversion $\hat{\xb}_0^{\text{raw}} = \hat{\xb}_0 \cdot \text{scale\_factor}$
must happen before calling \texttt{distribution\_aware\_utility()}, and the gradient
must flow through this scaling.

% ============================================================================
\section{DDIM Acceleration}
\label{sec:ddim}
% ============================================================================

Full DDPM requires $T - 1 = 999$ sequential UNet evaluations. DDIM (Song et al., 2021)
allows using a \textbf{subsequence} of timesteps with a deterministic update rule.

\subsection{DDIM update rule}

Choose a subsequence of $S$ timesteps: $\tau_1 < \tau_2 < \cdots < \tau_S$ from $\{1, \ldots, T\}$.
For example, with $S = 50$: $\tau = \{20, 40, 60, \ldots, 1000\}$ (uniform spacing).

At each DDIM step, going from $\tau_i$ to $\tau_{i-1}$ (denoting these as $t$ and $t'$ for clarity):
\begin{equation}
    \xb_{t'} = \sqrt{\alphabar_{t'}}\, \hat{\xb}_0
    + \sqrt{1 - \alphabar_{t'} - \sigma^2}\,
    \underbrace{\frac{\xb_t - \sqrt{\alphabar_t}\, \hat{\xb}_0}{\sqrt{1 - \alphabar_t}}}_{\text{``predicted noise direction''}}
    + \sigma\, \zb
    \label{eq:ddim}
\end{equation}

where $\hat{\xb}_0 = (\xb_t - \sqrt{1-\alphabar_t}\,\eb_\theta(\xb_t, t)) / \sqrt{\alphabar_t}$
and $\sigma$ controls stochasticity:
\begin{itemize}[nosep]
    \item $\sigma = 0$: deterministic DDIM (fully reproducible given $\xb_T$)
    \item $\sigma = \sqrt{(1-\alphabar_{t'})/(1-\alphabar_t)} \cdot \sqrt{\beta_t}$: recovers DDPM
\end{itemize}

\subsection{Guided DDIM}

For guided generation, replace $\hat{\xb}_0$ with $\hat{\xb}_0^{\text{guided}}$ from
Eq.~\eqref{eq:x0_guided}:
\begin{equation}
    \hat{\xb}_0^{\text{guided}}
    = \hat{\xb}_0 + \frac{w\,(1-\alphabar_t)}{\sqrt{\alphabar_t}}\, \gb_t
\end{equation}

Then use $\hat{\xb}_0^{\text{guided}}$ in the DDIM update (Eq.~\eqref{eq:ddim}).

Alternatively, one can compute the guided noise prediction $\hat{\eb}_{\text{guided}}$
(Eq.~\eqref{eq:eps_guided}) and derive $\hat{\xb}_0^{\text{guided}}$ from it. Both are equivalent.

\begin{algorithm}[H]
\caption{Guided DDIM (Deterministic, $\sigma = 0$)}
\label{alg:guided_ddim}
\begin{algorithmic}[1]
\Require UNet $\eb_\theta$, schedule, GP model, guidance scale $w$,
         number of steps $S$, firing rate cap $f_{\max}$
\Ensure Generated image $\xb_0$ with high utility
\State Compute timestep subsequence $\tau_S > \tau_{S-1} > \cdots > \tau_1$
    \Comment{e.g., uniform in $[1, T{-}1]$}
\State $\xb_{\tau_S} \sim \mathcal{N}(\mathbf{0}, \Ib)$
\For{$i = S, S{-}1, \ldots, 1$}
    \State $t \gets \tau_i$, \quad $t' \gets \tau_{i-1}$ (or $0$ if $i = 1$)
    \Statex
    \State \textbf{--- Compute Tweedie estimate with gradient tracking ---}
    \State $\xb_t.\text{requires\_grad} \gets \text{True}$
    \State $\hat{\eb} \gets \eb_\theta(\xb_t, t)$
    \State $\hat{\xb}_0 \gets (\xb_t - \sqrt{1-\alphabar_t}\,\hat{\eb}) / \sqrt{\alphabar_t}$
    \Statex
    \State \textbf{--- Utility guidance ---}
    \State $U, \mu_g \gets \texttt{utility}(\hat{\xb}_0)$ in raw pixel space
    \If{$\exp(\mu_g) < f_{\max}$}
        \State $\gb_t \gets \nabla_{\xb_t} U(\hat{\xb}_0)$ via \texttt{backward()}
    \Else
        \State $\gb_t \gets \mathbf{0}$
    \EndIf
    \Statex
    \State \textbf{--- Guided Tweedie ---}
    \State $\hat{\xb}_0^g \gets \hat{\xb}_0 + \frac{w(1-\alphabar_t)}{\sqrt{\alphabar_t}}\,\gb_t$
    \Statex
    \State \textbf{--- DDIM step ---}
    \State \textbf{direction} $\gets (\xb_t - \sqrt{\alphabar_t}\,\hat{\xb}_0^g) / \sqrt{1-\alphabar_t}$
    \If{$i > 1$}
        \State $\xb_{t'} \gets \sqrt{\alphabar_{t'}}\,\hat{\xb}_0^g
            + \sqrt{1-\alphabar_{t'}}\cdot\textbf{direction}$
    \Else
        \State $\xb_0 \gets \hat{\xb}_0^g$
            \Comment{Final step: just return the clean estimate}
    \EndIf
    \State Detach $\xb_{t'}$ from computation graph
\EndFor
\State \Return $\xb_0$
\end{algorithmic}
\end{algorithm}

\subsection{Choosing the number of DDIM steps $S$}

For a first implementation, $S = 50$ is a good starting point.
Quality degrades below $S \approx 20$ for typical diffusion models.

With $S = 50$ steps, each requiring one UNet forward + backward + one utility evaluation:
\begin{itemize}[nosep]
    \item UNet forward: $\sim$5\,ms (2.16M params, 64$\times$64, GPU)
    \item UNet backward: $\sim$10\,ms (frozen weights, only input gradient)
    \item Utility evaluation: $\sim$1\,ms (\texttt{distribution\_aware\_utility}, single image)
    \item Utility backward: $\sim$2\,ms (gradient through kernel)
    \item \textbf{Total per step}: $\sim$18\,ms
    \item \textbf{Total for 50 steps}: $\sim$0.9\,s
\end{itemize}

This is very feasible. Even full DDPM at 999 steps would take $\sim$18\,s per image.

% ============================================================================
\section{Gradient Computation: Full vs.\ Approximate}
\label{sec:gradient_approx}
% ============================================================================

The guidance gradient $\gb_t = \nabla_{\xb_t} U(\hat{\xb}_0(\xb_t, t))$ involves the
chain rule through the Tweedie formula and the UNet:
\begin{equation}
    \gb_t = \underbrace{\nabla_{\hat{\xb}_0} U}_{\text{utility gradient}}
    \cdot
    \underbrace{\frac{\partial \hat{\xb}_0}{\partial \xb_t}}_{\text{Tweedie Jacobian}}
\end{equation}

\subsection{Full gradient (through UNet)}

From Eq.~\eqref{eq:tweedie_jacobian}:
\begin{equation}
    \frac{\partial \hat{\xb}_0}{\partial \xb_t}
    = \frac{1}{\sqrt{\alphabar_t}}
    \left(\Ib - \sqrt{1 - \alphabar_t}\, \frac{\partial \eb_\theta}{\partial \xb_t}\right)
\end{equation}

In PyTorch, this is computed automatically: we make $\xb_t$ a leaf tensor with gradients,
run the UNet, compute Tweedie, evaluate utility, and call \texttt{backward()}.
The UNet's parameters are frozen (\texttt{requires\_grad=False}), but the computation
graph through the UNet's \emph{forward pass} is retained for the input gradient.

\textbf{Cost}: One UNet forward pass (already needed for denoising) + one UNet backward
pass (for the gradient). Roughly $2\times$ the cost of a standard unconditional step.

\subsection{Approximate gradient (skip UNet Jacobian)}
\label{sec:approx_grad}

An often-used approximation ignores the UNet's Jacobian entirely:
\begin{equation}
    \gb_t^{\text{approx}}
    = \frac{1}{\sqrt{\alphabar_t}}\, \nabla_{\hat{\xb}_0} U(\hat{\xb}_0)
    \label{eq:g_approx}
\end{equation}

This treats $\hat{\xb}_0$ as if it depends on $\xb_t$ only through the direct
$\xb_t / \sqrt{\alphabar_t}$ term, ignoring the $\eb_\theta(\xb_t, t)$ dependence.

\textbf{Justification}: The UNet Jacobian $\partial \eb_\theta / \partial \xb_t$
is a complicated nonlinear function. In many practical settings (text-to-image, etc.),
the approximate gradient works nearly as well as the full gradient. The intuition:
the UNet's prediction $\eb_\theta$ is a ``local correction'' whose Jacobian is
close to zero-mean (it doesn't systematically align with or oppose the utility gradient).

\textbf{Advantage}: No backward pass through the UNet. The cost per step drops from
$\sim$18\,ms to $\sim$8\,ms. For 50 DDIM steps: $\sim$0.4\,s vs.\ $\sim$0.9\,s.

\textbf{When this approximation breaks down}: When the UNet's correction is large
(early steps, high noise) or when the utility gradient is very sensitive to the
specific noise prediction. At low noise levels ($t$ small), the approximation is
excellent because $\sqrt{1-\alphabar_t}$ is small, making the UNet term negligible.

\subsection{Recommendation for implementation}

\textbf{Start with the approximate gradient} (Eq.~\eqref{eq:g_approx}):
\begin{enumerate}[nosep]
    \item Simpler implementation (no gradient tracking through UNet)
    \item Faster per step (no UNet backward)
    \item Good enough for most of the reverse process (especially late steps)
\end{enumerate}

Then implement the full gradient as a comparison. If the approximate gradient
produces images with comparable utility and naturalness, keep it.

\textbf{Hybrid approach}: Use the approximate gradient for early steps ($t > 200$)
where the Tweedie estimate is rough anyway, and the full gradient for late steps ($t < 200$)
where precision matters. This gives most of the benefit with minimal extra cost.

% ============================================================================
\section{Practical Considerations}
\label{sec:practical}
% ============================================================================

\subsection{Guidance scale $w$}

The guidance scale $w$ controls the tradeoff between naturalness and utility.
$w = 0$ gives unconditional sampling (natural but random utility).
Large $w$ pushes hard toward high utility but can distort the image.

\textbf{Calibration}: At early steps ($t$ large), $\hat{\xb}_0$ is a rough estimate
and utility gradients are noisy. At late steps ($t$ small), they are precise.
The factor $(1-\alphabar_t)/\sqrt{\alphabar_t}$ in Eq.~\eqref{eq:x0_guided} naturally
downweights guidance at low noise levels. This is desirable: most of the ``utility steering''
should happen at intermediate noise levels where the image structure is being decided.

Starting point: $w \in [0.1, 10]$. Print the utility of $\hat{\xb}_0$ and the guidance
gradient norm at each step to calibrate.

\subsection{Firing rate guard}

At early steps, $\hat{\xb}_0$ may have extreme pixel values (the Tweedie estimate from
heavily noisy images is rough). These can produce unrealistic firing rates.
The $f_{\max}$ guard in the algorithm zeros out the guidance gradient when
$\exp(\mu_g) > f_{\max}$. This prevents the guidance from chasing
physiologically impossible utility values at early steps.

\subsection{RF mask and image reconstruction}

The utility depends only on RF pixels ($\sim$21\% of the 64$\times$64 image).
But the reverse process generates \emph{all} pixels. The guidance gradient
$\gb_t$ will be nonzero only at RF pixels (because the utility gradient is zero
at non-RF pixels, since the kernel masks them out). The diffusion model shapes
the entire image for naturalness; the utility guides only the RF region.

This is an \textbf{advantage} over Approach A: the non-RF region is naturally
generated (not fixed at a starting image). The full image will look natural.

\subsection{Multiple samples}

Each run with a different $\xb_T$ seed produces a different image. To find the
best one:
\begin{enumerate}[nosep]
    \item Generate $K$ images (e.g., $K = 10$) with different seeds
    \item Evaluate $U(\xb_0)$ on each
    \item Select the one with highest utility
\end{enumerate}

With $\sim$1\,s per image, generating 10 candidates takes $\sim$10\,s. This is
comparable to the 50-step SGD optimization in Approach A.

\subsection{No starting point control}

Unlike Approach A (which starts from a chosen image), Approach D starts from noise.
The generated image will be ``natural with high utility'' but not a modification of
any specific image. This is a fundamental difference in the problem formulation:
\begin{itemize}[nosep]
    \item Approach A: ``make this image more informative while keeping it natural''
    \item Approach D: ``generate a natural image that is maximally informative''
\end{itemize}

If starting from a specific image is needed, \textbf{SDEdit} (Meng et al., 2022)
offers a middle ground: encode the starting image to an intermediate noise level $t^*$
(add noise via the forward process), then run guided reverse from $t^*$ to 0.
This preserves the coarse structure of the starting image while refining details
under guidance. This is a natural extension but not part of the first implementation.

% ============================================================================
\section{Comparison with Approach A}
\label{sec:comparison}
% ============================================================================

\begin{center}
\begin{tabular}{lcc}
\toprule
\textbf{Property} & \textbf{A: L2 Tweedie} & \textbf{D: Guided Reverse} \\
\midrule
Optimization target & Modify existing image & Generate from noise \\
UNet usage & 1 fwd per step (detached) & 1 fwd+bwd per step \\
Score evaluation & At fixed $t^*$ & At correct $t$ \\
In-distribution? & Sometimes (input can drift) & Always (reverse process) \\
Signal at $\mathcal{M}$ & Weak (Tweedie $\approx$ identity) & Strong (score at all scales) \\
Norm control & No & Yes (generation is bounded) \\
\# UNet calls & 50 fwd & 50 fwd + 50 bwd (DDIM-50) \\
Wall time & $\sim$3\,s & $\sim$1\,s (no LBFGS inner loop) \\
Starting image & User-chosen & Random (or SDEdit) \\
Output diversity & Deterministic (fixed start) & Stochastic (different seeds) \\
\bottomrule
\end{tabular}
\end{center}

A key advantage of Approach D: the diffusion model is \emph{always} evaluated
at the noise level it was trained for. In Approach A, after a few optimization steps,
the image $\xc$ may have amplified norms or broken correlations, making it
out-of-distribution for the UNet at any fixed $t^*$.

% ============================================================================
\section{Summary of What to Implement}
\label{sec:implementation}
% ============================================================================

\subsection{Core function: \texttt{guided\_reverse\_diffusion()}}

\begin{enumerate}
    \item \textbf{Inputs}: trained UNet, schedule, GP model + likelihood, guidance scale $w$,
          number of DDIM steps $S$, $f_{\max}$, optional seed.
    \item \textbf{Build DDIM timestep subsequence}: $S$ timesteps uniformly spaced in $[1, T{-}1]$.
    \item \textbf{Initialize}: $\xb_{\tau_S} \sim \mathcal{N}(\mathbf{0}, \Ib)$ on device.
    \item \textbf{Loop} over DDIM steps (Section~\ref{sec:ddim}, Algorithm~\ref{alg:guided_ddim}).
    \item \textbf{Return}: final $\xb_0$ in raw pixel space + per-step metrics
          (utility of $\hat{\xb}_0$, guidance gradient norm, $\alphabar_t$).
\end{enumerate}

\subsection{Supporting changes}

\begin{itemize}[nosep]
    \item \texttt{diffusion\_model.py}: add \texttt{ddim\_step()} function
          (non-\texttt{@torch.no\_grad}, takes $\xb_t$, $t$, $t'$, noise prediction).
    \item \texttt{guided\_optimization.py}: add \texttt{guided\_reverse\_diffusion()},
          modify \texttt{main()} to call it as an additional optimization path alongside
          utility-only and Approach A.
    \item Plotting: show generated image, per-step utility trajectory, compare with
          Approach A result.
\end{itemize}

\subsection{Validation plan}

\begin{enumerate}[nosep]
    \item \textbf{$w = 0$}: should match unconditional sampling quality.
    \item \textbf{Small $w$}: utility should increase modestly vs.\ unconditional.
    \item \textbf{Large $w$}: utility should be high; check if image still looks natural.
    \item \textbf{Compare}: utility and naturalness of Approach D output vs.\ Approach A output.
    \item \textbf{Multiple seeds}: verify that best-of-$K$ selection finds high-utility images.
\end{enumerate}

\end{document}
